\scanpage{10}
\deleted{avec des conditions (1.5) (1.6) qu'on n'explicite pas.}

Prenons le cas (3.12) $B=B'\times B''$, $f=(f',f'')$, $g(X',X'')=g'(X')\times g''(X'')$, \deleted{alors les données sont, si on pose concernant $\alpha$} supposons \unsure{que|de plus} $f'g'\simeq\mathrm{id}_{B^{\unsure{''|'}}}$, \deleted{\unsure{$f'g'$}} $f''g''\simeq\mathrm{id}_{B''}$ (i.e. $g',g''$ pl. fidèles), donc $\varphi$ est connu quand on connaît
\[
 \boxed{\varphi'':B'\longrightarrow B'',\qquad \varphi':B''\mathrel{\unsure{-|\to}}B'} \tag*{(3.13)}
\]
\deleted{et la donnée des $\alpha$ \unsure{revient à} $\lambda_{kji}$ revient à celle des}
\[
 \deleted{X'\times\varphi'(\deleted{\unsure{B}}X'')\longrightarrow X'} \tag*{\deleted{(3.10)}}
\]
\[
 \deleted{\varphi''(X')\times X''\longrightarrow X''} \tag*{\deleted{(3.11)}}
\]
\[
 \deleted{\left\{\begin{aligned}
 X'&\longrightarrow\varphi'\varphi''(X')\qquad X''\mathrel{\unsure{-|\to}}\varphi''\varphi'(X'')\\
 &X
 \end{aligned}\right.}
\]
\deleted{Il y a huit $\lambda_{kji}$ ($2^3$), dont deux triviaux ($i=j=k$), donc reste 6 non triviaux $\lambda_{1,1,2}:$} Le cas a) ci-dessus ne peut se présenter, et il reste donc des données
\[
 \boxed{\left\{\begin{aligned}
 \mathrm{id}_{B'}&\xrightarrow{\lambda'}\varphi'\varphi''\\
 \mathrm{id}_{B''}&\xrightarrow{\lambda''}\varphi''\varphi'
 \end{aligned}\right.} \tag*{(3.14)}
\]
Il faudrait encore \deleted{vérifier} expliciter dans ce cas les conditions (1.5) et (1.6) --

\diagram{\resizebox{\linewidth}{!}{\begin{tikzcd}[ampersand replacement=\&,column sep=small]
 \varphi(X',X'')\arrow[r,"\lambda"] \& \varphi^2(X',X'')\arrow[r,shift left,"\text{sont égaux ?}"]\arrow[r,shift right] \& \varphi(X',X'')\\[-8pt]
 \shortparallel \& \shortparallel \& \shortparallel\\[-8pt]
 (X'\times\varphi'(X''),X''\times\varphi''(X'))\arrow[r] \&
 \begin{gathered}\bigl(X'\times\varphi'(X'')\times\varphi'(X'')\times\varphi'\varphi''X',\\[-2pt]
 X''\times\varphi''(X')\times\varphi''(X')\times\varphi''\varphi'(X'')\bigr)\end{gathered} \&
 \begin{gathered}(X'\times\varphi'(X''),\\ X''\times\varphi''(X'))\end{gathered}
\end{tikzcd}}}
\margin{cf. 1.5 semble O.K.}
\deleted{Par contre, il semble que 1.6 impose des restrictions à $\lambda',\lambda''$}
Il \deleted{\unsure{ne}} semble \deleted{pas} qu'il n'y ait \deleted{pas} les conditions \ldots
\[
 \boxed{\begin{aligned}
 \varphi''*\lambda'&=\deleted{\unsure{\varphi}}\lambda''*\varphi'':\varphi''\longrightarrow\varphi''\varphi'\varphi''\,;\\
 \lambda'*\varphi'&=\varphi'*\lambda'':\varphi'\deleted{\unsure{\varphi}}\longrightarrow\varphi'\varphi''\varphi'
 \end{aligned}}
\]
Explicitons la catégorie $A$ en termes de $\varphi',\varphi'',\lambda',\lambda''$

Elle est formée des couples $(X',\unsure{\alpha''|X''},u',u'')$
\[
 X'\xrightarrow{u'}\varphi'(X''),\qquad X''\xrightarrow{u''}\varphi''(X') \tag*{(3.15)}
\]
\deleted{$X'\xrightarrow{\unsure{u'}}$} définissant
% archive sheet number: 9
